\documentclass[11pt,a4paper]{article}
\usepackage[utf8]{inputenc}
\usepackage[T1]{fontenc}
\usepackage[french]{babel}
\usepackage{amsmath,amsfonts,amssymb,amsthm}
\usepackage{geometry}
\geometry{margin=1in}
\usepackage{listings}
\usepackage{newunicodechar}

\newunicodechar{ℕ}{\mathbb{N}}
\newunicodechar{∀}{\forall}
\newunicodechar{∃}{\exists}
\newunicodechar{→}{\rightarrow}
\newunicodechar{∈}{\in}
\newunicodechar{ℝ}{\mathbb{R}}
\newunicodechar{≥}{\ge}
\newunicodechar{∧}{\land}

\title{Bornes Rigoureuses sur le Nombre de Tur\'an pour les Cycles Pairs}
\author{Charles EDOU NZE, chercheur ind\'ependant}
\date{}

\begin{document}
\maketitle

\begin{abstract}
Nous pr\'esentons une analyse structurale rigoureuse du nombre de Tur\'an pour les cycles pairs. Charles EDOU NZE, chercheur ind\'ependant.
\end{abstract}

\section{Cadre Axiomatique}
Soit $G = (V, E)$ un graphe. Le nombre de Tur\'an $\mathrm{ex}(n, H)$ est le nombre maximal d'ar\^etes dans un graphe \`a $n$ sommets ne contenant pas $H$.

\section{Lemmes Structuraux}
\subsection{Lemme 1 : Sous-graphes Bipartis de Haut Degr\'e}
Si $G$ a un degr\'e moyen $d$, il existe un sous-graphe biparti o\`u chaque sommet d'une partition a un degr\'e d'au moins $d/4$.
\begin{proof}
Soit une partition al\'eatoire. L'esp\'erance des ar\^etes est $|E|/2$. L'\'elimination it\'erative des sommets de degr\'e inf\'erieur \`a $d/4$ pr\'eserve un sous-graphe non vide.
\end{proof}

\section{Architecture d'Autoformalisation}
\begin{lstlisting}[basicstyle=\ttfamily\small, breaklines=true, mathescape=true]
import Mathlib.Combinatorics.SimpleGraph.Basic

variable {V : Type*} [Fintype V]
variable (G : SimpleGraph V)

theorem bipartite_high_degree (d : $\mathbb{R}$) (h : G.avg_degree $\ge$ d) :
  $\exists$ (A B : Set V) (G' : SimpleGraph V), G'.IsBipartite $\land$ $\forall$ v $\in$ B, G'.degree v $\ge$ d / 4 := by
  admit
\end{lstlisting}

\end{document}
